% Figure: order of 2 mod 7 — powers cycle 1,2,4 with period 3.
% Source for Figure~\ref{fig:orderwheel}. TikZ only; regenerable by compiling main.tex.
\begin{tikzpicture}[
  el/.style={circle,draw,minimum size=7mm,font=\footnotesize},
  lab/.style={font=\footnotesize\sffamily,text=slate,align=center}
]
\foreach \k/\v in {0/1,1/2,2/4}
  \node[el] (n\k) at (90-120*\k:1.5cm) {\v};
\draw[->,thick,accent] (n0) to[bend left=18] node[midway,above right,lab] {$\cdot 2$} (n1);
\draw[->,thick,accent] (n1) to[bend left=18] (n2);
\draw[->,thick,accent] (n2) to[bend left=18] node[midway,left,lab] {$2^3 = 8 \equiv 1$} (n0);
\node[lab] at (0,-2.6) {$\operatorname{ord}_7(2) = 3$: the loop
$1 \to 2 \to 4 \to 1$;\\
$2^k$ depends only on $k \bmod 3$, and $3 \mid 7-1$.};
\node[lab] at (4.6,0.9) {Compare: $\operatorname{ord}_7(3) = 6$};
\node[el] (m0) at (4.6,0.15) {1};
\node[lab,align=center]{$3^1,3^2,\dots$ hits all\\
six nonzero residues before};
\node[lab,align=center]{the order always divides\\
$\varphi(7) = 6$};
\end{tikzpicture}
